\section{Finite polygamma evaluation and a digamma master}
\label{sec:integer}

The one-index reduction is useful only if its new kernel can be handled
explicitly. At integer orders it is finite; at confluence its apparent
singularities disappear; and the entire positive-order family is
generated by one function of two shifts.

\subsection{Partial fractions with all simple poles paired}

\begin{proposition}\label{prop:partial-fractions}
Let $p,q$ be positive integers and $\Rea x,\Rea y>0$. If $x\ne y$, put
$D=y-x$ and
\begin{align}
 A_k&=(-1)^{p-k}\binom{p+q-k-1}{q-1}D^{-(p+q-k)},\quad 1\le k\le p,
 \label{eq:pf-a}\\
 B_l&=(-1)^p\binom{p+q-l-1}{p-1}D^{-(p+q-l)},\quad 1\le l\le q.
 \label{eq:pf-b}
\end{align}
Then $A_1+B_1=0$ and
\begin{equation}\label{eq:finite-kernel}
 \Kern(p,q;x,y)=A_1\{\psi(y)-\psi(x)\}
 +\sum_{k=2}^{p}A_k\zeta(k,x)
 +\sum_{l=2}^{q}B_l\zeta(l,y).
\end{equation}
The empty sums are zero. At $x=y$ the value is $\zeta(p+q,x)$.
\end{proposition}

\begin{proof}
Expand $(X+D)^{-q}$ at $X=0$, and $X^{-p}$ at $X=-D$.
The respective principal parts of $X^{-p}(X+D)^{-q}$ have coefficients
\eqref{eq:pf-a} and \eqref{eq:pf-b}. Their difference from that rational
function is entire and vanishes at infinity, hence is zero.
Its $X^{-1}$ coefficient at infinity gives $A_1+B_1=0$.
Consequently the only two series which would diverge separately occur
as the convergent pair
\[
 A_1\sum_{n\ge0}\left(\frac1{n+x}-\frac1{n+y}\right)
 =A_1\{\psi(y)-\psi(x)\}.
\]
All higher powers sum absolutely. This proves \eqref{eq:finite-kernel}.
The confluent value follows directly from the original kernel.
\end{proof}

More generally, confluence comes with every Taylor coefficient:
\begin{equation}\label{eq:kernel-confluence}
 \Kern(s,t;x,x+h)=\sum_{r=0}^{\infty}
 \frac{(-1)^r(t)_r}{r!}h^r\zeta(s+t+r,x),\qquad |h|<|x|.
\end{equation}
Here $\Rea x>0$ and $\Rea(x+h)>0$; on compact sets a uniformly
smaller disk is understood.
The identity follows first for $\Rea(s+t)>1$ by the binomial theorem,
since $|n+x|\ge |x|$, and then meromorphically. It is preferable to
subtracting nearly cancelling negative powers of $D$ when computing
close shifts. The finite formula and \eqref{eq:classical-ladders}
establish a finite polygamma evaluation of every
$\Ord(p,q;a,b)$ with $p\ge1$, $q\ge2$.

\subsection{The exact parity left by equal shifts}

\begin{theorem}\label{thm:parity}
For $p\ge1$, $q\ge2$, $N=p+q$ and $\Rea a>0$, define
\[
 L_j=\begin{cases}\binom{N-j-1}{q-1},&j\le p,\\0,&j>p,\end{cases}
 \qquad
 R_j=\begin{cases}\binom{N-j-1}{p-1},&j\le q,\\0,&j>q.\end{cases}
\]
Then
\begin{align}
 \Ord(p,q;a,a)={}&E_p(a)\zeta(q,a+1)\notag\\
 &+\frac{(-1)^q}{4}\sum_{j=2}^{\max(p,q)}
 \{1-(-1)^{N-j}\}\{L_j+(-1)^jR_j\}\zeta(j,a+1)\notag\\
 &-\frac14\sum_{k=1}^{p}\binom{N-k-1}{q-1}
 \{(-1)^q(a+1)^{-k}+(-1)^{p-k}a^{-k}\}.
 \label{eq:ordered-parity}
\end{align}
Only orders $j$ with $p+q+j$ odd occur in the correction in the second
line. When $p=q$, that entire line vanishes.
\end{theorem}

\begin{proof}
Apply \eqref{eq:finite-kernel} to
$\Kern(p,q;a+2,a+1)-\Kern(p,q;a,a+1)$, so that $D=-1$ and $D=1$.
Use $\zeta(j,a)=a^{-j}+\zeta(j,a+1)$ and
$\zeta(j,a+2)=\zeta(j,a+1)-(a+1)^{-j}$, together with the same
shift recurrence for $\psi$. For $j\ge2$ the resulting coefficient is
$(-1)^q\{1-(-1)^{N-j}\}\{L_j+(-1)^jR_j\}/4$.
The simple-pole pairs contribute only the finite third line.
If $p=q$, then $L_j=R_j$: the second brace vanishes for odd $j$,
and the first vanishes for even $j$.
\end{proof}

This is a statement about explicit coefficients. It makes no claim
that the special values remaining in a particular formula are
arithmetically independent.

\begin{example}[The ordering survives]\label{ex:ordered-values}
Two equal-shift evaluations are
\begin{align}
 \Ord(2,3;1,1)&=\frac38\zeta(3)-\frac32\zeta(2)+\frac{33}{16},
 \label{eq:example-23}\\
 \Ord(3,2;1,1)&=\frac{31}{16}\zeta(2)-\frac{95}{32}.
 \label{eq:example-32}
\end{align}
Their average is the previously polarized example
$7\zeta(2)/32+3\zeta(3)/16-29/64$ in \cite{repo:finite}.
Confluent shifts give a different cancellation:
\begin{align}
 \Ord(2,2;1,2)&=\frac14\zeta(4)+\frac14\zeta(2)-\frac9{16},
 \label{eq:example-1122}\\
 \Ord(2,2;2,1)&=\frac7{36}\zeta(2)-\frac14\zeta(4)-\frac{11}{576}.
 \label{eq:example-2211}
\end{align}
The $\zeta(4)$ term is an exact confluence contribution erased when
the two assignments are added.
\end{example}

\subsection{One master function generates all positive orders}

Define the holomorphic divided difference
\[
 \Psi[x,y]=\begin{cases}
 \dfrac{\psi(y)-\psi(x)}{y-x},&x\ne y,\\[5pt]
 \psi'(x),&x=y.
 \end{cases}
\]
For $\Rea a,\Rea b>0$ set
\begin{equation}\label{eq:digamma-master}
 F(a,b)=-\frac{\psi(b+1)}{2a(a+1)}
 +\frac14\{\Psi[a+2,b+1]-\Psi[a,b+1]\}.
\end{equation}

\begin{theorem}[Digamma master]\label{thm:digamma-master}
For $p\ge1$ and $q\ge2$,
\begin{equation}\label{eq:master-derivatives}
 \Ord(p,q;a,b)=
 \frac{(-1)^{p+q-2}}{(p-1)!(q-1)!}
 \partial_a^{p-1}\partial_b^{q-1}F(a,b).
\end{equation}
For real $a,b>0$ and $|x|<a$, $|y|<b$ this is equivalently
\begin{equation}\label{eq:master-geometric}
 \sum_{p\ge1}\sum_{q\ge2}x^{p-1}y^{q-1}\Ord(p,q;a,b)
 =F(a-x,b-y)-F(a-x,b).
\end{equation}
Both apparent singularities $b=a\pm1$ in a quotient version of $F$
are removable.
\end{theorem}

\begin{proof}
The residue of $\Ord(1,t;a,b)$ at $t=1$ is
$E_1(a)=1/(2a(a+1))$. The constant of the Hurwitz zeta factor there is
$-\psi(b+1)$, while $\Kern(1,1;x,y)=\Psi[x,y]$ is convergent.
Thus \eqref{eq:digamma-master} is precisely
$\FP_{t=1}\Ord(1,t;a,b)$ with $a,b$ fixed.

Direct differentiation of the weights gives
\begin{equation}\label{eq:weight-derivatives}
 \partial_a^k w_n(s,a)=(-1)^k(s)_k w_n(s+k,a).
\end{equation}
For $q\ge2$,
\[
 \partial_b^{q-1}\Ord(1,t;a,b)
 =(-1)^{q-1}(t)_{q-1}\Ord(1,t+q-1;a,b).
\]
The $b$ derivative kills the residue, and the right side is regular at
$t=1$. Its value is $(-1)^{q-1}(q-1)!\Ord(1,q;a,b)$.
Differentiate in $a$ and use \eqref{eq:weight-derivatives} again.
This proves \eqref{eq:master-derivatives}; Taylor expansion proves
\eqref{eq:master-geometric}.

There is also a direct convergence interpretation. Summing the
geometric weight expansions produces $w_m(1,a-x)$ in the inner
index and $w_n(1,b-y)-w_n(1,b)$ in the outer index. The latter
difference has one additional power of decay. Their original ordered
sum is absolutely convergent on the stated polydisk. This proves that
\eqref{eq:master-geometric} is an ordinary generating identity.
\end{proof}

The slightly larger domain $a>0$, $b>-1$, $|x|<a$, $|y|<b+1$
also works: the outer summation starts at $n=1$. We retain the
smaller common shift domain elsewhere to keep the original definitions
uniform.

\subsection{Nonlinear harmonic weights and Gamma ratios}

Let $\Bell_r$ denote the complete exponential Bell polynomial, normalized by
\begin{equation}\label{eq:bell-definition}
 \exp\left(\sum_{j\ge1}x_j\frac{v^j}{j!}\right)
 =\sum_{r\ge0}\Bell_r(x_1,\ldots,x_r)\frac{v^r}{r!}.
\end{equation}
For $h_j(n,a)=\sum_{\ell=0}^{2}(n+a+\ell)^{-j}$, one has
\begin{equation}\label{eq:bell-weight}
 w_n(p,a)=\frac{\Gamma(n+a)}{\Gamma(n+a+3)}
 \frac{\Bell_{p-1}(h_1,1!h_2,\ldots,(p-2)!h_{p-1})}{(p-1)!}.
\end{equation}
For $p=1$ the Bell polynomial is $\Bell_0=1$.
Indeed $w_n(1,a)=\Gamma(n+a)/\Gamma(n+a+3)$, and the logarithmic
derivatives of this ratio are $(-1)^j(j-1)!h_j$.
Equation \eqref{eq:weight-derivatives} and the Bell formula for
successive derivatives prove \eqref{eq:bell-weight}.

This translates the ordered identities into explicit nonlinear
harmonic sums. For example, write
$P_n=(n+1)(n+2)(n+3)$ and
$h_{j,n}=H_{n+3}^{(j)}-H_n^{(j)}$. Then \eqref{eq:example-23} becomes
\begin{equation}\label{eq:harmonic-example}
 \sum_{0\le m<n}\frac{(n-m)^2h_{1,m}(h_{1,n}^2+h_{2,n})}{P_mP_n}
 =\frac34\zeta(3)-3\zeta(2)+\frac{33}{8}.
\end{equation}
All higher Bell combinations are supplied by the same finite master,
with no new summation coordinates. At integral shifts they involve
finite harmonic differences; at arbitrary shifts the $h_j$ are
finite polygamma differences by the usual recurrence.

\subsection{A Gamma primitive under common translation}

For $a,b>0$, $h\ge0$, define
$\Lambda_h(x)=\log\Gamma(x+h)-\log\Gamma(x)$.
Then
\begin{equation}\label{eq:gamma-primitive}
 \begin{split}
 \int_0^h\Def(1,1;a+c,b+c)\dd c
 ={}&\frac{\Lambda_h(b+1)-\Lambda_h(a+2)}{b-a-1}\\
 &-\frac{\Lambda_h(b+1)-\Lambda_h(a)}{b-a+1}.
 \end{split}
\end{equation}
At a zero denominator its quotient is the confluent derivative
$\Lambda_h'(x)=\psi(x+h)-\psi(x)$. To prove the identity, replace
each $\Kern(1,1)$ by its digamma divided difference; common
translation leaves its denominator constant, and the primitive of
$\psi$ is $\log\Gamma$. This also proves the holomorphic path version
in the right half-plane. For positive integral $h$, the Gamma ratios
are finite products, so the primitive becomes a finite logarithmic
identity. Further argument derivatives generate polygamma primitives
of all positive-order kernels.
